% Luc Destombe --- licence CC BY-NC-SA 4.0
% https://creativecommons.org/licenses/by-nc-sa/4.0/deed.fr
% Fichier autonome : compiler avec pdflatex, deux fois (signets, nombre de pages).
\documentclass[10pt,a4paper,twocolumn]{article}
\makeatletter
\newif\iflycee@unecolonne
\newif\iflycee@palatino
\lycee@palatinotrue

\RequirePackage[utf8]{inputenc}
\RequirePackage[T1]{fontenc}
\RequirePackage[french]{babel}
\RequirePackage{etoolbox}

\newcommand{\ShorthandsOff}{\@ifundefined{shorthandoff}{}{\shorthandoff{;:!?}}}
\newcommand{\ShorthandsOn}{\@ifundefined{shorthandon}{}{\shorthandon{;:!?}}}
\ShorthandsOff
\AtBeginDocument{\ShorthandsOn}
\AtBeginEnvironment{tikzpicture}{\ShorthandsOff}

\RequirePackage{geometry}
\RequirePackage{fancyhdr}
\RequirePackage{lastpage}
\RequirePackage{multicol}
\RequirePackage{array}
\RequirePackage{enumitem}
\RequirePackage{xcolor}
\RequirePackage{needspace}

\RequirePackage{amsmath,amssymb}
\RequirePackage{mathpazo}
\DeclareMathSizes{10.95}{10.95}{8}{5.5}
\begingroup\catcode`\_=\active \catcode`\^=\active
\gdef_#1{\sb{\scriptscriptstyle #1}}
\gdef^#1{\sp{\scriptscriptstyle\lycee@moinsepais #1}}
\endgroup
\newcommand\lycee@moins{\mathbin{%
  \setbox\z@\hbox{$\m@th\scriptscriptstyle\lycee@moinsorig$}%
  \dimen@=.55\wd\z@ \advance\dimen@-.5pt
  \hbox to.75\wd\z@{\hss$\m@th\scriptscriptstyle\vcenter{\hbox{%
    \tikz\draw[line width=.5pt,line cap=round](0,0)--(\the\dimen@,0);}}$\hss}}}
\begingroup\lccode`\~=`\- \lowercase{\endgroup\let~}\lycee@moins
\newcommand\lycee@moinsepais{\mathcode`\-="8000 }
\AtBeginDocument{\mathchardef\lycee@moinsorig=\mathcode`\-\relax}
\AtBeginDocument{\catcode`\_=12 \mathcode`\_="8000
  \catcode`\^=12 \mathcode`\^="8000 }
\newcommand\lycee@exposants{%
  \fontdimen13\textfont2=.48\fontdimen6\textfont2
  \fontdimen14\textfont2=.48\fontdimen6\textfont2
  \fontdimen15\textfont2=.48\fontdimen6\textfont2
  \fontdimen13\scriptfont2=.48\fontdimen6\scriptfont2
  \fontdimen14\scriptfont2=.48\fontdimen6\scriptfont2
  \fontdimen15\scriptfont2=.48\fontdimen6\scriptfont2}
\every@math@size\expandafter{\the\every@math@size\lycee@exposants}
\newlength{\retraitcalculs}
\setlength{\retraitcalculs}{2em}

\RequirePackage{tikz}
\usetikzlibrary{arrows.meta,calc,babel,shapes.misc}
\RequirePackage[most]{tcolorbox}
\tcbuselibrary{skins,breakable}

\definecolor{coulDef}{HTML}{1F4E79}
\definecolor{coulProp}{HTML}{7030A0}
\definecolor{coulRem}{HTML}{C55A11}
\definecolor{coulEx}{HTML}{2E7D32}
\definecolor{coulExo}{HTML}{B03A2E}
\definecolor{coulPart}{HTML}{1F4E79}
\definecolor{coulMeth}{HTML}{1B6E4F}
\definecolor{coulCourbe}{HTML}{0B5ED7}
\definecolor{coulCourbeB}{HTML}{C00000}
\definecolor{coulCourbeC}{HTML}{2E7D32}
\definecolor{coulGrille}{HTML}{8C8C8C}
\definecolor{coulRep}{HTML}{1B6E4F}
\definecolor{coulBar}{HTML}{2E7D32}

\newcommand{\R}{\mathbb{R}}
\newcommand{\Rs}{\mathbb{R}^{*}}
\renewcommand{\le}{\leqslant}
\renewcommand{\ge}{\geqslant}
\renewcommand{\approx}{\simeq}
\newcommand{\vect}[1]{\overrightarrow{#1}}
\newcommand{\norme}[1]{\left\lVert #1 \right\rVert}
\newcommand{\intervalle}[2]{\left[#1\,;#2\right]}
\RequirePackage{cancel}
\renewcommand{\CancelColor}{\color{coulExo}}
\newcommand{\fois}[1]{\times{\color{coulExo}#1}}
\newcommand{\barre}[1]{\cancel{#1}}

\tikzset{grille/.style={coulGrille,line width=0.4pt}}
\newcommand{\quadrillage}[4]{\draw[grille,step=1] (#1,#2) grid (#3,#4);}
\tikzset{
  croix/.style={cross out,draw,line width=0.6pt,inner sep=0pt,minimum size=4pt},
  croixdroite/.style={croix,rotate=45,line width=0.8pt},
}
\newcommand{\pt}[3]{\path (#1) node[croix]{} node[#2,font=\small,inner sep=2.5pt]{$#3$};}

\newcommand{\lycee@repere}[4]{%
  \draw[grille,step=1] (#1,#2) grid (#3,#4);
  \draw[-{Stealth[length=2mm]},thick] (#1,0)--({#3+0.4},0) node[below left=-1pt]{$x$};
  \draw[-{Stealth[length=2mm]},thick] (0,#2)--(0,{#4+0.4}) node[below left=-1pt]{$y$};
  \node[below left,font=\scriptsize,inner sep=1.5pt] at (0,0) {$0$};}
\newcommand{\lycee@gradx}[1]{\foreach \k/\l in {#1}{%
  \draw (\k,0.07)--(\k,-0.07) node[below,font=\scriptsize,inner sep=1.5pt]{$\l$};}}
\newcommand{\lycee@grady}[1]{\foreach \k/\l in {#1}{%
  \draw (0.07,\k)--(-0.07,\k) node[left,font=\scriptsize,inner sep=1.5pt]{$\l$};}}
\newcommand{\lycee@intff}[2]{\left[#1\,;#2\right]}
\newcommand{\lycee@intfo}[2]{\left[#1\,;#2\right[}
\newcommand{\lycee@intof}[2]{\left]#1\,;#2\right]}
\newcommand{\lycee@intoo}[2]{\left]#1\,;#2\right[}
\newcommand{\lycee@ens}[1]{\left\{#1\right\}}
\AtBeginDocument{%
  \providecommand{\repere}{\lycee@repere}%
  \providecommand{\gradx}{\lycee@gradx}%
  \providecommand{\grady}{\lycee@grady}%
  \providecommand{\Cf}{\mathcal{C}\sb{\scriptscriptstyle f}}%
  \providecommand{\Cg}{\mathcal{C}\sb{\scriptscriptstyle g}}%
  \providecommand{\intff}{\lycee@intff}%
  \providecommand{\intfo}{\lycee@intfo}%
  \providecommand{\intof}{\lycee@intof}%
  \providecommand{\intoo}{\lycee@intoo}%
  \providecommand{\ens}{\lycee@ens}}

\newlist{questions}{enumerate}{3}
\setlist[questions,1]{label=\arabic*\textdegree),leftmargin=*,itemsep=2pt,topsep=3pt}
\setlist[questions,2]{label=\alph*),leftmargin=*,itemsep=1pt,topsep=2pt}
\setlist[questions,3]{label=\roman*),leftmargin=*,itemsep=1pt,topsep=1pt}
\setlist[itemize]{label=\textbullet,leftmargin=*,itemsep=2pt,topsep=3pt}

\newcommand{\besoinplace}[1]{\par\penalty\@M\begingroup
  \dimen@=#1\relax\dimen@ii=\pagegoal\advance\dimen@ii-\pagetotal
  \ifdim\dimen@>\dimen@ii\ifdim\dimen@ii>\z@\vfil\fi\break\fi\endgroup}

\newcommand{\lycee@pied}{\thepage/\pageref{LastPage}}
\newcommand{\entete}[2]{%
  \pagestyle{fancy}\fancyhf{}%
  \fancyhead[L]{\footnotesize\color{coulPart}\textbf{#1}}%
  \fancyhead[R]{\footnotesize\color{coulPart}\textbf{#2}}%
  \fancyfoot[C]{\footnotesize\lycee@pied}%
  \renewcommand{\headrulewidth}{0.4pt}%
  \renewcommand{\headrule}{\hbox to\headwidth{%
    \color{coulPart!40}\leaders\hrule height \headrulewidth\hfill}}}

\RequirePackage{ccicons}
\newcommand{\lycee@licence}{{\scriptsize\color{gray}%
  \copyright\ Luc Destombe --- \ccbyncsa\ CC BY-NC-SA 4.0}}
\newif\iflycee@licence \lycee@licencetrue
\newcommand{\sanslicence}{\lycee@licencefalse}
\AtBeginDocument{\iflycee@licence\AddToHookNext{shipout/foreground}{%
  \put(0,0){\rlap{\kern\dimexpr1in+\hoffset+\oddsidemargin+\textwidth\relax
    \llap{\raisebox{-\dimexpr1in+\voffset+\topmargin+\headheight+\headsep
      +\textheight+\footskip\relax}{\lycee@licence}}}}}\fi}

\newcommand{\formulesserrees}{%
  \setlength{\abovedisplayskip}{3pt}\setlength{\belowdisplayskip}{3pt}%
  \setlength{\abovedisplayshortskip}{2pt}\setlength{\belowdisplayshortskip}{3pt}}

\setlength{\parindent}{0pt}

\RequirePackage[implicit=false,hidelinks,pdfencoding=auto,psdextra,bookmarksopen,
  bookmarksopenlevel=1,pdfcreator={},pdfproducer={}]{hyperref}
\RequirePackage{bookmark}
\hypersetup{pdfauthor={Luc Destombe},pdfkeywords={CC BY-NC-SA 4.0}}
\newcounter{lycee@signet}
\newcommand{\signet}[3][]{\stepcounter{lycee@signet}%
  \edef\lycee@dest{\if\relax\detokenize{#1}\relax
    signet.\arabic{lycee@signet}\else#1\fi}%
  \hypertarget{\lycee@dest}{}\bookmark[level=#2,dest=\lycee@dest]{#3}}
\pdfstringdefDisableCommands{%
  \def\R{R}\def\vect#1{#1}\def\Cf{Cf}\def\Cg{Cg}\def\fois#1{×#1}%
  \def\barre#1{#1}\def\intervalle#1#2{[#1;#2]}\def\intff#1#2{[#1;#2]}%
  \def\textsuperscript#1{#1}\def\dfrac#1#2{#1/#2}\def\frac#1#2{#1/#2}%
  \def\vec#1{#1⃗}}%
\RequirePackage[official]{eurosym}

\tcbset{exercice corrige/.style={
  enhanced,breakable,before upper=\raggedright,
  colback=white,colframe=coulExo,
  boxrule=0.8pt,arc=2pt,
  left=6pt,right=6pt,top=8pt,bottom=5pt,
  before skip=9pt,after skip=4pt,
  coltitle=white,fonttitle=\bfseries\small,
  attach boxed title to top left={xshift=8pt,yshift=-\tcboxedtitleheight/2},
  boxed title style={colback=coulExo,arc=2pt,boxrule=0pt}}}
\newcounter{exo}
\newtcolorbox{exercice}[1][]{exercice corrige,
  phantom={\signet[exercice-\arabic{exo}]{2}{Exercice \arabic{exo}}},
  title={Exercice \theexo\ifx\relax#1\relax\else{} --- #1\fi}}
\newenvironment{exo}[1][\relax]{\refstepcounter{exo}\begin{exercice}[#1]}{\end{exercice}}

\RequirePackage{cuted}
\geometry{top=1.6cm,bottom=1cm,left=1cm,right=1cm,headheight=14pt,footskip=0.6cm}

\newcommand{\entetecorrige}[3]{%
  \begin{strip}
  \begin{center}
    \begin{tcolorbox}[enhanced,width=0.92\linewidth,colback=white,colframe=coulPart,
        boxrule=1.6pt,arc=3pt,halign=center,top=5pt,bottom=5pt]
      {\LARGE\bfseries\color{coulPart} #1}\\[3pt]
      {\normalsize #2}
    \end{tcolorbox}
  \end{center}
  \if\relax\detokenize{#3}\relax\else

  \vspace{2pt}
  \noindent
  \begin{tcolorbox}[enhanced,colback=white,colframe=black!35,boxrule=0.5pt,arc=2pt,
      left=7pt,right=7pt,top=4pt,bottom=4pt]
  {\small #3}
  \end{tcolorbox}
  \vspace{6pt}
  \fi
  \end{strip}}

\newcounter{partie}
\newcommand{\partiebandeau}[1]{%
  \besoinplace{8\baselineskip}\vspace{6pt}%
  \begin{tcolorbox}[enhanced,colback=coulPart!8,colframe=coulPart,boxrule=0pt,
      leftrule=3pt,arc=0pt,left=5pt,right=4pt,top=3pt,bottom=3pt,
      before skip=4pt,after skip=2pt]
    \bfseries\color{coulPart}#1\signet{1}{#1}
  \end{tcolorbox}}
\newcommand{\partie}[1]{\stepcounter{partie}\partiebandeau{\Roman{partie}) #1}}
\newcommand{\partiesansnum}[1]{\partiebandeau{#1}}

\newcommand{\res}[1]{%
  \tcbox[on line,colback=coulPart!7,colframe=coulPart,boxrule=0.5pt,arc=1pt,
    boxsep=0pt,left=3pt,right=3pt,top=1pt,bottom=2pt]{$#1$}}
\newcommand{\calc}[1]{\par\smallskip\noindent$\displaystyle #1$\par}
\newenvironment{calculs}{\par\smallskip\noindent$\displaystyle\begin{aligned}[t]}%
  {\end{aligned}$\par\smallskip}
\newcommand{\rem}[1]{\par\smallskip{\small\itshape\color{coulPart}#1}\par}
\newcommand{\deux}[2]{\par\noindent\makebox[0.5\linewidth][l]{#1}#2}

\setlist[questions,1]{label=\arabic*\textdegree),leftmargin=*,itemsep=2pt,topsep=2pt}
\setlist[questions,2]{label=\alph*),leftmargin=*,itemsep=2pt,topsep=1pt}
\newlist{reponses}{enumerate}{1}
\setlist[reponses]{label=\alph*),leftmargin=*,itemsep=2pt,topsep=2pt}

\setlength{\parskip}{1pt}
\setlength{\columnsep}{0.6cm}
\raggedbottom
\formulesserrees
\AtBeginDocument{\formulesserrees}

\makeatother
\usepackage{tkz-tab}
\providecommand{\justif}[1]{\quad\text{\small\itshape\color{coulPart}#1}}
\entete{Ch 3 --- Calcul littéral et second degré --- Corrigé des exercices}{1\textsuperscript{re} mathématiques spécifiques}

\newcommand{\varmin}[3]{\par\smallskip{\centering
  \begin{tikzpicture}[scale=0.8]
  \tkzTabInit[lgt=1.1,espcl=1.5]{$x$/0.7,$#1$/1.1}{,$#2$,}
  \tkzTabVar{+/ ,-/$#3$,+/ }
  \end{tikzpicture}\par}}
\newcommand{\varmax}[3]{\par\smallskip{\centering
  \begin{tikzpicture}[scale=0.8]
  \tkzTabInit[lgt=1.1,espcl=1.5]{$x$/0.7,$#1$/1.1}{,$#2$,}
  \tkzTabVar{-/ ,+/$#3$,-/ }
  \end{tikzpicture}\par}}
\newcommand{\varint}[3]{\par\smallskip{\centering
  \begin{tikzpicture}[scale=0.8]
  \tkzTabInit[lgt=1.1,espcl=1.6]{$x$/0.7,$#1$/1.1}{#2}
  \tkzTabVar{#3}
  \end{tikzpicture}\par}}

\begin{document}

\entetecorrige{CALCUL LITTÉRAL ET SECOND DEGRÉ}{Corrigé --- Ch 3 : Fiche d'exercices --- 1\textsuperscript{re} mathématiques spécifiques}
{Les remarques en italique signalent les erreurs les plus fréquentes et les méthodes à
retenir. Les résultats sont encadrés. Tous les calculs se font sans calculatrice.}

\partie{Développer}

\begin{exo}[Distributivité simple]
\deux{$A=\res{5x+15}$}{$B=\res{8x-28}$}
\deux{$C=\res{-3x-18}$}{$D=-8+10x=\res{10x-8}$}
\deux{$E=\res{21-7x}$}{$F=\res{-x+9}$}
\rem{Le facteur multiplie \emph{chaque} terme : $-2\times(-5x)=+10x$.}
\end{exo}

\begin{exo}[Un facteur en $x$]
\deux{$A=\res{x^{2}+8x}$}{$B=\res{6x^{2}-2x}$}
\deux{$C=\res{-x^{2}+4x}$}{$D=\res{-8x^{2}-20x}$}
\deux{$E=\res{15x-6x^{2}}$}{}
\rem{$x\times x=x^{2}$ et $2x\times 3x=6x^{2}$.}
\end{exo}

\begin{exo}[Double distributivité]
\begin{calculs}
A &= x^{2}+5x+2x+10 &&\justif{quatre produits}\\
A &= \res{x^{2}+7x+10}
\end{calculs}
\deux{$B=\res{x^{2}+x-12}$}{$C=\res{x^{2}-7x+6}$}
\deux{$D=\res{2x^{2}+7x+3}$}{$E=\res{3x^{2}+4x-4}$}
\begin{calculs}
F &= 4x-x^{2}+20-5x &&\justif{$5\times(-x)=-5x$}\\
F &= \res{-x^{2}-x+20}
\end{calculs}
\end{exo}

\begin{exo}[Vrai ou faux ?]
\begin{reponses}
  \item Faux : $3(x+4)=\res{3x+12}$.
  \item \res{\text{Vrai}}.
  \item \res{\text{Vrai}}.
  \item Faux : $(x+1)(x+2)=\res{x^{2}+3x+2}$.
  \item Faux : $-(3-x)=\res{-3+x}$.
\end{reponses}
\rem{d) On oublie souvent les deux produits « croisés » $x\times 2$ et $1\times x$.}
\end{exo}

\partie{Identités remarquables}

\begin{exo}[Le carré d'une somme]
\begin{calculs}
A &= x^{2}+2\times x\times 4+4^{2} &&\justif{$a=x$, $b=4$}\\
A &= \res{x^{2}+8x+16}
\end{calculs}
\deux{$B=\res{x^{2}+2x+1}$}{$C=\res{4x^{2}+12x+9}$}
\deux{$D=\res{9x^{2}+30x+25}$}{}
\rem{Le double produit ne s'oublie pas : $(x+4)^{2}\ne x^{2}+16$.}
\end{exo}

\begin{exo}[Le carré d'une différence]
\begin{calculs}
A &= x^{2}-2\times x\times 3+3^{2} &&\justif{$a=x$, $b=3$}\\
A &= \res{x^{2}-6x+9}
\end{calculs}
\deux{$B=\res{x^{2}-20x+100}$}{$C=\res{4x^{2}-4x+1}$}
\deux{$D=\res{25x^{2}-20x+4}$}{}
\end{exo}

\begin{exo}[Somme par différence]
\deux{$A=x^{2}-6^{2}=\res{x^{2}-36}$}{$B=\res{x^{2}-1}$}
\deux{$C=(3x)^{2}-2^{2}=\res{9x^{2}-4}$}{$D=\res{16x^{2}-25}$}
\rem{$(3x)^{2}=9x^{2}$ : le $3$ est aussi élevé au carré.}
\end{exo}

\begin{exo}[La bonne réponse]
\deux{1. \res{\text{b}}\; ($x^{2}-14x+49$)}{2. \res{\text{c}}\; ($9x^{2}+6x+1$)}
\deux{3. \res{\text{b}}\; ($(x-4)(x+4)$)}{4. \res{\text{a}}\; ($4x^{2}-25$)}
\end{exo}

\partie{Équations du premier degré}

\begin{exo}[Premières équations]
\deux{a) $x=12-7$ : \res{S=\ens{5}}}{b) $x=-20\div 4$ : \res{S=\ens{-5}}}
\deux{c) $3x=15$ : \res{S=\ens{5}}}{}
\begin{calculs}
5x+2 &= 3x+10\\
2x &= 8\\
x &= 4
\end{calculs}
d) \res{S=\ens{4}} ; \quad e) $6x=12$ : \res{S=\ens{2}} ; \quad f) $12=4x$ :
\res{S=\ens{3}}
\end{exo}

\begin{exo}[Avec des parenthèses]
\begin{calculs}
3(x+2) &= 21\\
3x+6 &= 21 &&\justif{on développe d'abord}\\
3x &= 15\\
x &= 5
\end{calculs}
a) \res{S=\ens{5}} ; \quad b) $2x-8=x+1$ : \res{S=\ens{9}}
\begin{calculs}
5(x+1) &= 2(x+7)\\
5x+5 &= 2x+14\\
3x &= 9\\
x &= 3
\end{calculs}
c) \res{S=\ens{3}} ; \quad d) $-x+6=2x$ soit $6=3x$ : \res{S=\ens{2}}\par
e) $8x-4=3x+21$ soit $5x=25$ : \res{S=\ens{5}}
\end{exo}

\begin{exo}[L'inconnue au dénominateur]
\begin{calculs}
\dfrac{12}{x} &= 3\\
12 &= 3x &&\justif{on multiplie par $x$}\\
x &= 4
\end{calculs}
a) \res{S=\ens{4}} ; \quad b) $30=6x$ : \res{S=\ens{5}}\par
c) $8=-2x$ : \res{S=\ens{-4}} ; \quad d) $21=7x$ : \res{S=\ens{3}}\par
e) $\dfrac{120}{x}=15$ soit $120=15x$ : \res{x=8} amis.
\rem{Contrôle : $120\div 8=15$.}
\end{exo}

\begin{exo}[Deux formules pour la piscine]
\begin{questions}
  \item Entrées seules : $6\times 10=\res{60\text{ €}}$ ; carte :
        $40+4\times 10=\res{80\text{ €}}$.
  \item Entrées seules : \res{6x} ; carte : \res{40+4x}.
  \item \begin{calculs}
        6x &= 40+4x\\
        2x &= 40\\
        x &= 20
        \end{calculs}
        Les deux formules coûtent autant pour \res{20} entrées ($120$~€).
  \item Entrées seules : $150$~€ ; carte : $140$~€. La \res{\text{carte}} est moins
        chère.
\end{questions}
\end{exo}

\partie{Équations $x^{2}=a$}

\begin{exo}[Résoudre $x^{2}=a$]
\deux{a) \res{S=\ens{-4\,;4}}}{b) \res{S=\ens{-1\,;1}}}
\deux{c) \res{S=\ens{0}}}{d) \res{S=\ens{-\sqrt{3}\,;\sqrt{3}}}}
\deux{e) \res{S=\varnothing}}{f) \res{S=\ens{-12\,;12}}}
\rem{e) Un carré n'est jamais négatif. Ne pas oublier la solution négative.}
\end{exo}

\begin{exo}[Isoler $x^{2}$ d'abord]
\begin{calculs}
2x^{2} &= 98\\
x^{2} &= 49 &&\justif{on divise par $2$}
\end{calculs}
\deux{a) $x^{2}=9$ : \res{S=\ens{-3\,;3}}}{b) $x^{2}=36$ : \res{S=\ens{-6\,;6}}}
\deux{c) \res{S=\ens{-7\,;7}}}{d) $x^{2}=4$ : \res{S=\ens{-2\,;2}}}
\deux{e) $x^{2}=-9$ : \res{S=\varnothing}}{f) $x^{2}=25$ : \res{S=\ens{-5\,;5}}}
\end{exo}

\begin{exo}[Des carrés]
\begin{questions}
  \item $x^{2}=81$, donc $x=-9$ ou $x=9$ ; une longueur est positive : \res{9\text{ m}}.
  \item $x^{2}=100$ : côté \res{10\text{ cm}}.
  \item \res{-7} et \res{7}.
\end{questions}
\end{exo}

\partie{Équations produit}

\begin{exo}[Produit nul]
\begin{calculs}
x-1 &= 0 \quad\text{ou}\quad x+7=0\\
x &= 1 \quad\text{ou}\quad x=-7
\end{calculs}
\deux{a) \res{S=\ens{-7\,;1}}}{b) \res{S=\ens{-5\,;-2}}}
\deux{c) \res{S=\ens{3\,;8}}}{d) \res{S=\ens{-4\,;0}}}
\end{exo}

\begin{exo}[Des coefficients devant $x$]
\begin{calculs}
2x-4 &= 0 \quad\text{ou}\quad x+3=0\\
x &= 2 \quad\text{ou}\quad x=-3
\end{calculs}
\deux{a) \res{S=\ens{-3\,;2}}}{b) \res{S=\ens{-3\,;5}}}
\deux{c) \res{S=\ens{-4\,;3}}}{d) \res{S=\ens{2\,;6}}}
\deux{e) \res{S=\ens{-1\,;5}}}{}
\rem{d) $-x+6=0$ donne $x=6$ ; e) $-2x=-10$ donne $x=5$.}
\end{exo}

\begin{exo}[Factoriser par $x$ d'abord]
\begin{calculs}
x^{2}-4x &= 0\\
x(x-4) &= 0\\
x &= 0 \quad\text{ou}\quad x=4
\end{calculs}
\deux{a) \res{S=\ens{0\,;4}}}{b) $x(x+6)=0$ : \res{S=\ens{-6\,;0}}}
\deux{c) $x(3x-9)=0$ : \res{S=\ens{0\,;3}}}{d) $x(2x+10)=0$ : \res{S=\ens{-5\,;0}}}
e) $x^{2}-7x=0$ soit $x(x-7)=0$ : \res{S=\ens{0\,;7}}
\rem{e) Diviser par $x$ ferait perdre la solution $0$.}
\end{exo}

\partie{Fonction polynôme du second degré}

\begin{exo}[Reconnaître, lire $a$, $b$ et $c$]
\deux{$f$ : \res{a=2\,;\ b=5\,;\ c=-1}}{$g$ : \res{a=-1\,;\ b=0\,;\ c=3}}
\deux{$h$ : \res{\text{pas du 2\textsuperscript{d} degré}}}{$k$ : \res{a=-1\,;\ b=4\,;\ c=0}}
\deux{$m$ : \res{a=1\,;\ b=0\,;\ c=0}}{$p$ : \res{a=3\,;\ b=-2\,;\ c=6}}
\rem{On lit chaque coefficient avec son signe, même si les termes sont dans le
désordre ($k$, $p$).}
\end{exo}

\begin{exo}[Calculer des images]
\begin{calculs}
f(-2) &= (-2)^{2}-3\times(-2)+2\\
f(-2) &= 4+6+2\\
f(-2) &= 12
\end{calculs}
1\textdegree) $f(0)=\res{2}$ ; $f(1)=\res{0}$ ; $f(-1)=\res{6}$ ; $f(3)=\res{2}$ ;
$f(-2)=\res{12}$.\par
2\textdegree) $g(0)=\res{5}$ ; $g(2)=-8+2+5=\res{-1}$ ; $g(-1)=-2-1+5=\res{2}$.
\rem{$(-1)^{2}=1$, donc $-2\times(-1)^{2}=-2$.}
\end{exo}

\begin{exo}[Sur la courbe ou non ?]
\deux{$f(1)=0$ : $A$ \res{\text{oui}}}{$f(-1)=-4$ : $B$ \res{\text{non}}}
\deux{$f(2)=5$ : $C$ \res{\text{oui}}}{$f(-3)=0$ : $D$ \res{\text{oui}}}
\deux{$f(0)=-3$ : $E$ \res{\text{non}}}{}
\end{exo}

\begin{exo}[Distance d'arrêt]
\begin{questions}
  \item $d(v)=v^{2}+3v$ a un terme en $v^{2}$ : \res{a=1\,;\ b=3\,;\ c=0}.
  \item $d(3)=\res{18}$ ; $d(5)=\res{40}$ ; $d(9)=\res{108}$. À $50$~km/h, la voiture
        s'arrête en $40$~m.
  \item $d(6)=36+18=\res{54}$ : la distance est multipliée par $3$, elle ne
        \res{\text{double pas}}.
\end{questions}
\end{exo}

\partie{Forme factorisée et racines}

\begin{exo}[Vérifier une forme factorisée]
\begin{calculs}
(x-2)(x+4) &= x^{2}+4x-2x-8\\
(x-2)(x+4) &= x^{2}+2x-8 &&\justif{on retrouve $f(x)$}\\[3pt]
3(x+2)(x-3) &= 3\left(x^{2}-3x+2x-6\right)\\
3(x+2)(x-3) &= 3x^{2}-3x-18\\[3pt]
-(x-1)(x-5) &= -\left(x^{2}-6x+5\right)\\
-(x-1)(x-5) &= -x^{2}+6x-5\\[3pt]
2(x-2)(x+2) &= 2\left(x^{2}-4\right)\\
2(x-2)(x+2) &= 2x^{2}-8
\end{calculs}
\rem{On part de la forme factorisée, la plus compliquée, et on retrouve l'autre.}
\end{exo}

\begin{exo}[Lire les racines]
\deux{$f$ : \res{3} et \res{-1}}{$g$ : \res{-5} et \res{2}}
\deux{$h$ : \res{4} et \res{6}}{$k$ : \res{0} et \res{7}}
\deux{$m$ : \res{-2} et \res{-8}}{}
\rem{Dans $(x+5)$, la racine est $-5$ : on change le signe.}
\end{exo}

\begin{exo}[Points sur les axes]
\begin{questions}
  \item $(x+2)(x-3)=x^{2}-3x+2x-6=x^{2}-x-6$.
  \item $x+2=0$ ou $x-3=0$ : racines \res{-2} et \res{3}.
  \item \res{(-2\,;0)} et \res{(3\,;0)}.
  \item $f(0)=-6$ : la courbe coupe l'axe des ordonnées en \res{(0\,;-6)}.
\end{questions}
\end{exo}

\partie{Représentation graphique}

\begin{exo}[Tableau de valeurs et tracé]
{\centering\small
\renewcommand{\arraystretch}{1.2}\setlength{\tabcolsep}{4pt}
\begin{tabular}{|c|*{7}{c|}}
\hline
$x$ & $-1$ & $0$ & $1$ & $2$ & $3$ & $4$ & $5$\\
\hline
$f(x)$ & $8$ & $3$ & $0$ & $-1$ & $0$ & $3$ & $8$\\
\hline
\end{tabular}\par\smallskip
\begin{tabular}{|c|*{7}{c|}}
\hline
$x$ & $-2$ & $-1$ & $0$ & $1$ & $2$ & $3$ & $4$\\
\hline
$g(x)$ & $-5$ & $0$ & $3$ & $4$ & $3$ & $0$ & $-5$\\
\hline
\end{tabular}\par}
\smallskip
\begin{minipage}[t]{0.48\linewidth}\centering
\begin{tikzpicture}[x=0.3cm,y=0.3cm]
  \repere{-2}{-2}{6}{9}
  \gradx{2,4} \grady{2,4,6,8}
  \draw[coulCourbe,line width=1pt,domain=-1:5,samples=60] plot (\x,{\x*\x-4*\x+3});
  \foreach \p in {(-1,8),(0,3),(1,0),(2,-1),(3,0),(4,3),(5,8)}
    \path[coulCourbe] \p node[croix]{};
\end{tikzpicture}\par
{\small a) $S(2\,;-1)$, axe $x=2$}
\end{minipage}\hfill
\begin{minipage}[t]{0.48\linewidth}\centering
\begin{tikzpicture}[x=0.3cm,y=0.3cm]
  \repere{-3}{-6}{5}{5}
  \gradx{-2,2,4} \grady{-4,-2,2,4}
  \draw[coulCourbe,line width=1pt,domain=-2:4,samples=60] plot (\x,{-\x*\x+2*\x+3});
  \foreach \p in {(-2,-5),(-1,0),(0,3),(1,4),(2,3),(3,0),(4,-5)}
    \path[coulCourbe] \p node[croix]{};
\end{tikzpicture}\par
{\small b) $S(1\,;4)$, axe $x=1$}
\end{minipage}
\rem{Les valeurs se répètent de part et d'autre du sommet : c'est un contrôle.}
\end{exo}

\begin{exo}[Associer une fonction à sa courbe]
$g$ et $k$ ont $a<0$ (courbes tournées vers le bas) : $g(0)=4$, donc
$g\to\res{\mathcal{C}_{1}}$ ; $k(0)=-1$, donc $k\to\res{\mathcal{C}_{4}}$.\par
$f$ et $h$ ont $a>0$ : $f(0)=-3$, donc $f\to\res{\mathcal{C}_{3}}$ ; $h(0)=0$, donc
$h\to\res{\mathcal{C}_{2}}$.
\rem{Deux indices suffisent : le signe de $a$ et le point $(0\,;c)$.}
\end{exo}

\begin{exo}[Lire une parabole]
\begin{questions}
  \item Vers le \res{\text{haut}} ; cohérent : $a=1>0$.
  \item Sommet \res{S(-1\,;-4)} ; axe : la droite verticale qui passe par \res{x=-1}.
  \item Racines \res{-3} et \res{1}.
  \item Points d'ordonnée $-3$ : \res{S=\ens{-2\,;0}}.
  \item $f(0)=\res{-3}$ : le point $(0\,;-3)$, sur l'axe des ordonnées.
\end{questions}
\end{exo}

\partie{Tableaux de variations}

\begin{exo}[Sommet et tableau de variations]
a) Racines $2$ et $6$ ; $x_{S}=\dfrac{2+6}{2}=4$ ; $f(4)=2\times(-2)=-4$ :
\res{S(4\,;-4)}, minimum ($a=1>0$).
\varmin{f(x)}{4}{-4}
b) Racines $-1$ et $3$ ; $x_{S}=1$ ; $g(1)=-2\times 2\times(-2)=8$ : \res{S(1\,;8)},
maximum ($a=-2<0$).
\varmax{g(x)}{1}{8}
c) Racines $-4$ et $-2$ ; $x_{S}=-3$ ; $h(-3)=3\times 1\times(-1)=-3$ :
\res{S(-3\,;-3)}, minimum.
\varmin{h(x)}{-3}{-3}
d) Racines $5$ et $-1$ ; $x_{S}=2$ ; $k(2)=-(-3)\times 3=9$ : \res{S(2\,;9)},
maximum.
\varmax{k(x)}{2}{9}
\rem{$a$ est le nombre devant les parenthèses : dans $-(x-5)(x+1)$, $a=-1$.}
\end{exo}

\begin{exo}[Factoriser d'abord]
a) $x(x-6)=0$ : racines $0$ et $6$ ; $x_{S}=3$ ; $f(3)=9-18=-9$ : \res{S(3\,;-9)},
minimum.
\varmin{f(x)}{3}{-9}
b) $x(-x+10)=0$ : racines $0$ et $10$ ; $x_{S}=5$ ; $g(5)=-25+50=25$ :
\res{S(5\,;25)}, maximum.
\varmax{g(x)}{5}{25}
c) $x(2x+8)=0$ : racines $0$ et $-4$ ; $x_{S}=-2$ ; $h(-2)=8-16=-8$ :
\res{S(-2\,;-8)}, minimum.
\varmin{h(x)}{-2}{-8}
d) $x^{2}=9$ : racines $-3$ et $3$ ; $x_{S}=0$ ; $k(0)=-9$ : \res{S(0\,;-9)},
minimum.
\varmin{k(x)}{0}{-9}
\end{exo}

\begin{exo}[Maximum ou minimum ?]
\begin{questions}
  \item $f$ : $a=4>0$, $x_{S}=2$, $f(2)=4\times 1\times(-1)$ : minimum \res{-4} en $2$.\par
        $g$ : $a=-1<0$, $x_{S}=-1$, $g(-1)=-2\times(-2)$ : maximum \res{4} en $-1$.\par
        $h$ : $a=-5<0$, $x_{S}=1$, $h(1)=-5\times 1\times(-1)$ : maximum \res{5} en $1$.
  \item $f(0)=4\times(-1)\times(-3)=\res{12}$ ; $f(4)=4\times 3\times 1=\res{12}$.
        \varint{f(x)}{$0$,$2$,$4$}{+/$12$,-/$-4$,+/$12$}
\end{questions}
\end{exo}

\partie{Problèmes type épreuve anticipée}

\begin{exo}[Un food truck]
\begin{questions}
  \item $B(0)=\res{-20}$ : sans aucun repas vendu, la journée coûte $200$~€.
  \item \begin{calculs}
        -(x-2)(x-10) &= -\left(x^{2}-12x+20\right)\\
        -(x-2)(x-10) &= -x^{2}+12x-20
        \end{calculs}
  \item $x=2$ ou $x=10$ : bénéfice nul pour \res{20} ou \res{100} repas.
  \item $x_{S}=6$ ; $B(6)=-(4)\times(-4)=16$ ; $B(12)=-(10)\times 2=-20$.
        \varint{B(x)}{$0$,$6$,$12$}{-/$-20$,+/$16$,-/$-20$}
  \item Bénéfice maximal \res{160\text{ €}} pour \res{60} repas.
\end{questions}
\rem{Penser aux unités : $x$ en dizaines de repas, $B(x)$ en dizaines d'euros.}
\end{exo}

\begin{exo}[Une balle lancée d'une falaise]
\begin{questions}
  \item $h(0)=\res{15}$ : la balle part de $15$~m, le haut de la falaise.
  \item \begin{calculs}
        -5(t+1)(t-3) &= -5\left(t^{2}-3t+t-3\right)\\
        -5(t+1)(t-3) &= -5\left(t^{2}-2t-3\right)\\
        -5(t+1)(t-3) &= -5t^{2}+10t+15
        \end{calculs}
  \item $t=-1$ ou $t=3$ ; un temps est positif : l'eau au bout de \res{3\text{ s}}.
  \item $t_{S}=\dfrac{-1+3}{2}=1$ ; $h(1)=-5\times 2\times(-2)=20$.
        \varint{h(t)}{$0$,$1$,$3$}{-/$15$,+/$20$,-/$0$}
  \item \res{20\text{ m}} au bout de \res{1\text{ s}}.
\end{questions}
\end{exo}

\begin{exo}[Le plus grand enclos]
\begin{questions}
  \item Deux côtés voisins font la moitié du périmètre, $20$~m : l'autre côté mesure
        \res{15\text{ m}}, l'aire \res{75\text{ m}^{2}}.
  \item De même, l'autre côté mesure $20-x$ et l'aire est $A(x)=x(20-x)$.
  \item $A(x)=\res{-x^{2}+20x}$ ; racines \res{0} et \res{20}.
  \item $x_{S}=10$ ; $A(10)=10\times 10=100$.
        \varint{A(x)}{$0$,$10$,$20$}{-/$0$,+/$100$,-/$0$}
  \item Un carré de \res{10\text{ m}} de côté : \res{100\text{ m}^{2}}.
\end{questions}
\end{exo}

\partiesansnum{Exercices type épreuve anticipée}

\begin{exo}[Automatismes (QCM)]
\deux{1. \res{\text{b}}\; ($4x=8$)}{2. \res{\text{b}}}
\deux{3. \res{\text{b}}}{4. \res{\text{b}}\; ($x=2$ ou $x=-4$)}
\deux{5. \res{\text{a}}\; ($4-6+1=-1$)}{6. \res{\text{b}}\; ($a<0$, milieu de $1$ et $5$)}
\end{exo}

\begin{exo}[Une épidémie de grippe]
\begin{questions}
  \item $f(2)=\res{12}$ : $1\,200$ malades à la semaine $2$.
  \item Pic à la semaine \res{4}, avec \res{1\,600} malades.
  \item $f(1)=-1+8=\res{7}$ : $700$ malades au bout d'une semaine.
  \item $-x(x-8)=-x^{2}+8x$. Racines $0$ et $8$ : l'épidémie s'arrête au bout de
        \res{8} semaines.
  \item $x_{S}=4$ ; $f(4)=-4\times(-4)=16$, soit \res{1\,600} malades.
        \varint{f(x)}{$0$,$4$,$8$}{-/$0$,+/$16$,-/$0$}
\end{questions}
\end{exo}

\begin{exo}[Le prix d'une place de théâtre]
\begin{questions}
  \item $200-80=\res{120}$ spectateurs ; recette $8\times 120=\res{960\text{ €}}$.
  \item \begin{calculs}
        200-10x &= 150\\
        -10x &= -50\\
        x &= 5
        \end{calculs}
        Une place à \res{5\text{ €}}.
  \item Recette = prix $\times$ spectateurs $=x(200-10x)$, puis
        $R(x)=200x-10x^{2}=-10x^{2}+200x$.
  \item $x=0$ ou $200-10x=0$ : racines \res{0} et \res{20}. Places gratuites ou
        plus aucun spectateur à $20$~€ : recette nulle.
  \item $x_{S}=10$ ; $R(10)=10\times 100=1\,000$.
        \varint{R(x)}{$0$,$10$,$20$}{-/$0$,+/$1\,000$,-/$0$}
  \item Place à \res{10\text{ €}} : recette maximale \res{1\,000\text{ €}}.
\end{questions}
\end{exo}

\end{document}
